\scanpage{56}
\mspage{12}

\deleted{\fbox{\unsure{4)}}} \inserted{(B)} \uline{Opération} \uline{de}\deleted{\uline{\unsure{affaiblissement|faiblissement} et de}} \inserted{sur}\uline{marquage}

\deleted{On peut la définir \deleted{par} \inserted{en termes de} la donnée, d'un pour chaque \deleted{arête du contour $R$} \inserted{composante connexe de} $\deleted{\unsure{R\times S}}\ \widetilde\Sigma\setminus R$ Soit $\deleted{\unsure{\widetilde\Sigma_*}}A_{\unsure{0|*}}^+$ eff. l'ens}

\deleted{\unsure{des a}} Soit $A^+$ l'ens des arêtes \inserted{(libres ou liées)} ``effectivement marquées'', \inserted{\unsure{de|du} $\Gamma$, \unsure{i.e.|p.ex.} telles que $r(a)>0$} \deleted{et $B$ ens des}, correspondant à des $C\in\pi_0(\Sigma)$ telles que $C\cap R\ne\varnothing$, soit $B^+$ l'ens des comp. connexes de $\Sigma^+\setminus R\cap\Sigma^+$ (lesquelles sont des int. ouverts), un surmarquage est défini par la donnée d'une application
\[
 B^+\xrightarrow{r'}\mathbb N,
\]

% slanted
\margin{\deleted{arêtes} cercles \deleted{\unsure{marqués}} \inserted{effectivement} marqués, \uline{arêtes} \inserted{\unsure{corr.}} sur $\Sigma$ (\uline{NB} elles se trouvent sur des cercles eff. marqués, et sont des intervalles ouverts)}

et peut se décrire \deleted{par} \inserted{comme} l'opération qui sur chaque \uline{arête} \inserted{(ouverte)} \deleted{qui} de $\Sigma$, correspondant à \unsure{un|une} $b\in B^+$, $r'(b)$ nouveaux sommets de \deleted{marquage} \inserted{rigidification}. Cette fois-ci, \deleted{\unsure{l'opéra}} il y a des choix arbitraires --- néanmoins l'opération donne un objet de $S\unsure{\mathscr T|\mathbb T}{\unsure{\mathbb D|D}}^{+}$, défini à isom unique près \inserted{soit $(X,\Sigma,R_1)$} --- muni bien entendu d'un ens marqué $R'\subset R_1$, tel que l'opération inverse de gommage de $R'$ redonne l'objet de départ $(X,\Sigma,R)$.
% archive sheet number: 55

